\chapter{Instalação Avançada}
\label{apendice:instalacao}

\section{Compilação com suporte a ODBC}

Para conectar \hayashi{} a bancos de dados via ODBC (SQL Server, Oracle, etc.):

\begin{lstlisting}[style=inline, language=sh]
cargo install hayashi-lang --features odbc
\end{lstlisting}

Requer driver ODBC instalado no sistema (\texttt{unixODBC} no Linux).

\section{Variáveis de ambiente}

\begin{center}
\begin{tabular}{ll}
\toprule
\textbf{Variável} & \textbf{Efeito} \\
\midrule
\texttt{HAY\_HISTORY}  & Caminho do arquivo de histórico do REPL \\
\texttt{HAY\_NOCOLOR}  & Desabilita cores nas mensagens de erro \\
\texttt{HAY\_TRACE}    & Habilita trace de execução (depuração) \\
\bottomrule
\end{tabular}
\end{center}

\section{Integração com editores}

\subsection*{Neovim}

Adicione ao \texttt{init.lua} para reconhecimento de sintaxe básico:

\begin{lstlisting}[style=inline, language=Hayashi]
vim.filetype.add({ extension = { hay = "hayashi" } })
\end{lstlisting}

Highlighting completo via Tree-sitter está no roadmap.

\subsection*{VS Code}

Extensão disponível na marketplace como \texttt{hayashi-lang}.
